%! Author = sriadityavedantam
%! Date = 3/25/26

\section{Functional Limits}

\lesson{28}{Wednesday 15 October 2025 9:10}{Day 28}

\begin{problem}{
    \[
        \lim_{x\to 3} x^2 = 9.
    \]
}
    Let $\eps > 0$ be given.
    We want to show $\abs{x^2 - 9} = \abs{x - 3}\abs{x + 3} < \eps$.
    Take $\delta \coloneqq \min\{1, \frac{\eps}{7}\}$.
    If $\abs{x-3}<\delta$, then $2 < x < 4$.
    Thus $\abs{x+3} < 7$.
    Hence
    \begin{align*}
        \abs{x^2 - 9} &= \abs{x-3}\abs{x+3}\\
        &< 7\abs{x-3}\\
        &< 7\delta\\
        &\leq \eps.
    \end{align*}
\end{problem}
\begin{theorem}{
    \[
        \lim_{x\to x_0} f(x) = L
    \]
    if and only if for every sequence $x_n\in A$, $x_n\neq x_0$, with $x_n\to x_0$, we have
    \[
        \lim_{n\to\infty} f(x_n) = L.
    \]
}
    \textbf{($\implies$).}
    Let $x_n\to x_0$ with $x_n\neq x_0$.
    Given $\eps > 0$, there exists $\delta > 0$ such that
    \[
        0 < \abs{x-x_0} < \delta \implies \abs{f(x)-L} < \eps.
    \]
    Since $x_n\to x_0$, there exists $N\in\n$ such that $n\geq N$ implies $\abs{x_n-x_0}<\delta$.
    Hence $\abs{f(x_n)-L}<\eps$.

    \textbf{($\impliedby$).}
    We prove the contrapositive.
    Suppose $\lim_{x\to x_0} f(x) \neq L$.
    Then there exists $\eps > 0$ such that for every $\delta > 0$, there exists $x\in A$, $x\neq x_0$, with
    \[
        0 < \abs{x-x_0} < \delta
    \]
    and
    \[
        \abs{f(x)-L} \geq \eps.
    \]
    For each $n\in\n$, choose $x_n\in A$ such that
    \[
        0 < \abs{x_n-x_0} < \frac{1}{n}.
    \]
    Then $x_n\to x_0$ and $\abs{f(x_n)-L} \geq \eps$.
    Hence $f(x_n)\not\to L$.
\end{theorem}

\begin{theorem}{
    Suppose $f,g : A\to\r$ and $c\in\r$.
    If
    \[
        \lim_{x\to c} f(x) = L,\quad \lim_{x\to c} g(x) = M,
    \]
    then:
    \begin{enumerate}
        \item $\lim_{x\to c} (f(x) + g(x)) = L + M$.
        \item $\lim_{x\to c} f(x)g(x) = LM$.
        \item $\lim_{x\to c} cf(x) = cL$.
        \item If $M\neq 0$, then $\lim_{x\to c} \dfrac{f(x)}{g(x)} = \dfrac{L}{M}$.
    \end{enumerate}
}
\end{theorem}

\begin{definition}[One-Sided Limits]
    Given $A\subset\r$, $c$ is a right limit point of $A$ if for all $\eps > 0$,
    \[
        (c, c+\eps)\cap A\neq\varnothing.
    \]
    Similarly, $c$ is a left limit point if for all $\eps > 0$,
    \[
        (c-\eps,c)\cap A\neq\varnothing.
    \]
\end{definition}

\begin{definition}[Right-Hand Limit]
    Suppose $c$ is a right limit point of $A$.
    We write
    \[
        \lim_{x\to c^+}f(x) = L
    \]
    if for every $\eps > 0$, there exists $\delta > 0$ such that
    \[
        c < x < c + \delta,\ x\in A \implies \abs{f(x) - L} < \eps.
    \]
\end{definition}

\begin{proposition}
    \[
        \lim_{x\to c^+}f(x) = L
    \]
    if and only if for every sequence $x_n\in A$ with $x_n > c$ and $x_n\to c$, we have
    \[
        \lim_{n\to\infty} f(x_n) = L.
    \]
\end{proposition}

\begin{proposition}
    Suppose $c$ is both a right and left limit point of $A$.
    Then
    \[
        \lim_{x\to c}f(x) = L \iff \lim_{x\to c^+} f(x)= \lim_{x\to c^-}f(x) = L.
    \]
\end{proposition}

\section{Continuous Functions}

\begin{definition}[Continuous]
    Suppose $A\subset\r$, $f:A\to \r$, and $c\in A$.
    We say $f$ is continuous at $c$ if for every $\eps > 0$, there exists $\delta > 0$ such that
    \[
        \abs{x - c} < \delta \implies \abs{f(x) - f(c)} < \eps.
    \]
\end{definition}

\begin{remark}
    Any function is continuous at any isolated point of $A$.
\end{remark}

\begin{definition}[One-Sided Continuity]
    $f$ is continuous from the right at $c$ if for every $\eps > 0$, there exists $\delta > 0$ such that
    \[
        0 < x-c < \delta \implies \abs{f(x) - f(c)}<\eps.
    \]
\end{definition}

\begin{example}
    \[
        f(x) = \begin{cases}
                   \dfrac{x}{\abs{x}}, &x\neq 0\\
                   1, &x = 0.
        \end{cases}
    \]
    Not left continuous at $0$, but right continuous at $0$.
\end{example}

\begin{example}
    \[
        S \coloneqq \left\{\dfrac{1}{n} : n\in\n\right\}.
    \]
    $0$ is a right limit point of $S$, but not a left limit point.
\end{example}

\begin{remark}
    If $c$ is a limit point of $A$ and $c\in A$, then $f:A\to\r$ is continuous at $c$ if and only if
    \[
        \lim_{x\to c} f(x) = f(c).
    \]
\end{remark}

\begin{example}
    \[
        f(x) = \lfloor x \rfloor.
    \]
    It is continuous at all $x\notin\z$.
    It is right continuous everywhere.
    It is not left continuous at integers.
\end{example}

\begin{example}
    \[
        f(x) = \begin{cases}
                   x\sin(1/x), &x\neq 0\\
                   0, &x = 0.
        \end{cases}
    \]
    Continuous for all $c\in\r$.
\end{example}

\begin{proposition}{
    This function is continuous at $0$.
}
    Let $\eps > 0$.
    If $\abs{x}<\eps$, then
    \[
        \abs{x\sin(1/x)} \leq \abs{x} < \eps.
    \]
\end{proposition}

\begin{theorem}{
    Suppose $f:A\to \r$, $f(A)\subset B$, and $g:B\to\r$.
    If $f$ is continuous at $c\in A$ and $g$ is continuous at $f(c)$, then $g\circ f$ is continuous at $c$.
}
    Let $\eps > 0$.
    Since $g$ is continuous at $f(c)$, there exists $\delta_1 > 0$ such that
    \[
        \abs{y - f(c)} < \delta_1 \implies \abs{g(y) - g(f(c))} < \eps.
    \]
    Since $f$ is continuous at $c$, there exists $\delta > 0$ such that
    \[
        \abs{x-c} < \delta \implies \abs{f(x)-f(c)} < \delta_1.
    \]
    Hence
    \[
        \abs{g(f(x)) - g(f(c))} < \eps.
    \]
\end{theorem}

\begin{definition}[Continuous]
    A function $f:A\to\r$ is continuous if it is continuous at every $c\in A$.
\end{definition}

\section{Continuous Functions on Compact Sets}

\begin{theorem}[Continuous Functions on Compact Sets]{
    Suppose $C\subset\r$ is compact and $f:C\to \r$ is continuous.
    Then $f(C)$ is compact.
}
    Let $y_n\in f(C)$ with $y_n = f(x_n)$ and $x_n\in C$.
    Since $C$ is compact, there exists a subsequence $x_{n_k}\to x\in C$.
    By continuity of $f$, we have
    \[
        f(x_{n_k}) \to f(x).
    \]
    Hence $y_{n_k}\to f(x)\in f(C)$.
\end{theorem}

\begin{corollary}
    $f$ is continuous at $c$ if and only if whenever $x_n\to c$ with $x_n\in A$, we have $f(x_n)\to f(c)$.
\end{corollary}